\documentclass[11pt,letterpaper]{article}

% Packages
\usepackage[margin=1in,headheight=14pt]{geometry}
\usepackage{fontenc}
\usepackage{inputenc}
\usepackage{hyperref}
\usepackage{booktabs}
\usepackage{tabularx}
\usepackage{array}
\usepackage{longtable}
\usepackage{xcolor}
\usepackage{titlesec}
\usepackage{parskip}
\usepackage{enumitem}
\usepackage{fancyhdr}
\usepackage[expansion=false]{microtype}

% Colors
\definecolor{accentblue}{RGB}{0, 82, 155}
\definecolor{lightgray}{RGB}{245, 245, 245}
\definecolor{darkgray}{RGB}{80, 80, 80}

% Hyperref setup
\hypersetup{
    colorlinks=true,
    linkcolor=accentblue,
    urlcolor=accentblue,
    pdftitle={EVN HyperLogistics --- RVCC Venture Competition Application},
    pdfauthor={Tony Nguyen, Daniel Evans, Joel Vinas}
}

% Section formatting
\titleformat{\section}{\large\bfseries\color{accentblue}}{}{0em}{}[\titlerule]
\titlespacing{\section}{0pt}{1.5em}{0.5em}

% Header / footer
\pagestyle{fancy}
\fancyhf{}
\renewcommand{\headrulewidth}{0.4pt}
\lhead{\textcolor{darkgray}{\small EVN HyperLogistics}}
\rhead{\textcolor{darkgray}{\small RVCC 2026 Submission}}
\cfoot{\textcolor{darkgray}{\small \thepage}}

% Table styling helpers
\newcolumntype{L}[1]{>{\raggedright\arraybackslash}p{#1}}
\newcolumntype{C}[1]{>{\centering\arraybackslash}p{#1}}

\begin{document}

% ─── Title Block ──────────────────────────────────────────────────────────────
\begin{center}
    {\LARGE\bfseries\color{accentblue} EVN HyperLogistics}\\[0.3em]
    {\large RVCC Venture Competition Application}\\[1em]
    \hrule height 1.5pt
    \vspace{0.8em}
    \begin{tabular}{@{}ll@{}}
        \textbf{Venture:} & EVN HyperLogistics \\[0.2em]
        \textbf{Product:} & HyperLogistics --- Safety-compliant, explainable AI rerouting for middle-mile logistics \\[0.2em]
        \textbf{Date:}    & March 15, 2026 \\[0.2em]
        \textbf{Team:}    & Tony Nguyen \textperiodcentered{} Daniel Evans \textperiodcentered{} Joel Vinas \\[0.2em]
        \textbf{Demo Video:} & \href{https://youtu.be/iZP4qxMR6w4}{https://youtu.be/iZP4qxMR6w4} \\
    \end{tabular}
    \vspace{0.8em}
    \hrule height 1.5pt
\end{center}

\vspace{1em}

% ─── Section 1 ────────────────────────────────────────────────────────────────
\section{Problem \& Value Proposition}
\textit{What problem are you solving and why does your solution improve the situation?}

\medskip
Middle-mile logistics managers face a critical daily problem: when disruptions hit --- weather events, traffic accidents, road closures --- they must make immediate rerouting decisions that are both operationally sound and physically compliant with DOT regulations. Existing tools can flag risk, but they stop short of producing a safe, explainable reroute recommendation.

We call this the \textbf{prediction-action gap}.

The scale of this problem is real. Our system is grounded in 7.7 million US accident records, multi-terabyte NOAA weather data, and 600{,}000+ DOT bridge constraint records. These are not model assumptions --- they define the actual operating environment freight teams work in every day.

\textbf{HyperLogistics closes this gap} by combining three components into a single decision flow:

\vspace{0.5em}
\begin{tabular}{@{}L{3.2cm}L{10.5cm}@{}}
    \toprule
    \textbf{Component} & \textbf{Role} \\
    \midrule
    \textbf{ReMindRAG} & Retrieves grounded evidence for reroute reasoning from historical disruption knowledge \\[0.4em]
    \textbf{SRSNet} & Forecasts short-horizon (4--8 hour) disruption propagation aligned to transit windows \\[0.4em]
    \textbf{Symbolic Safety Veto} & Hard-blocks any route that violates DOT bridge clearance or load limits \\
    \bottomrule
\end{tabular}
\vspace{0.5em}

The result: a dispatch recommendation that is not just intelligent, but \textbf{compliance-ready and explainable}. Every output cites its reasoning and has passed a deterministic safety check --- something no general-purpose AI tool currently provides for this workflow.

% ─── Section 2 ────────────────────────────────────────────────────────────────
\section{Customer \& Market Opportunity}
\textit{Who is the target customer and why is this opportunity meaningful?}

\medskip
\textbf{Primary users:} Logistics Network Managers and Area Managers at trucking carriers and third-party logistics providers (3PLs).

These operators own middle-mile execution outcomes. They are accountable for on-time performance, cost efficiency, and safety compliance under constant disruption pressure. Their core job during a disruption is to triage, reroute, and confirm DOT compliance --- often in minutes.

\textbf{Why this opportunity is meaningful:}

\begin{itemize}[leftmargin=1.5em, itemsep=0.2em]
    \item Disruptions are high-frequency and consequential. A single bad reroute can violate bridge load limits, incur DOT liability, or delay a full shipment lane.
    \item Most routing AI tools today optimize for efficiency but \textbf{lack deterministic safety enforcement}. They are not built for compliance-critical decisions.
    \item HyperLogistics sits at the exact intersection where this is unaddressed: \textbf{operations + explainability + hard compliance enforcement}.
\end{itemize}

To validate product-market fit, we built 100 domain-adapted instruction scenarios drawn directly from real manager decision patterns --- disruption queries, bridge constraint vetoes, and reroute approval flows. These are not synthetic edge cases. They reflect the actual language and decisions logistics managers face daily.

\textbf{Market opportunity:}

\vspace{0.5em}
\begin{tabular}{@{}L{1.8cm}L{4.5cm}C{2.2cm}L{4.5cm}@{}}
    \toprule
    \textbf{Layer} & \textbf{Market} & \textbf{2025 Size} & \textbf{Growth} \\
    \midrule
    \textbf{TAM} & Middle-Mile Logistics & \$110B & 7\% CAGR $\to$ \$218B by 2035 \\[0.3em]
    \textbf{SAM} & Route Optimization Software & \$7.9B & 13.3\% CAGR $\to$ \$16.8B by 2031 \\[0.3em]
    \textbf{SOM} & Explainable AI in Logistics & \$1.7B & 33.7\% CAGR $\to$ \$31.7B by 2035 \\
    \bottomrule
\end{tabular}
\vspace{0.5em}

EVN targets the SAM/SOM intersection --- the fastest-growing segment (33.7\% CAGR) within an active purchasing market (\$7.9B) where compliance-enforcing, explainable AI does not yet exist as a category product.

\vspace{0.3em}
\textit{\small Sources: ResearchAndMarkets / 360iResearch (TAM), Mordor Intelligence (SAM), Future Market Insights (SOM)}

% ─── Section 3 ────────────────────────────────────────────────────────────────
\section{Progress \& Learning}
\textit{What steps have you taken so far and what have you learned?}

\subsection*{What We Built}

\vspace{0.3em}
\begin{tabular}{@{}L{4.5cm}L{9cm}@{}}
    \toprule
    \textbf{Milestone} & \textbf{Evidence} \\
    \midrule
    \textbf{Live deployment} & HyperLogistics Streamlit app is publicly accessible at \href{https://cs5542hyperlogistics.streamlit.app/}{cs5542hyperlogistics.streamlit.app} \\[0.4em]
    \textbf{Snowflake-native pipeline} & Medallion architecture (Bronze/Silver/Gold) with Snowpipe, External Tables, and Spatial SQL \\[0.4em]
    \textbf{Domain adaptation (Week 8)} & Accuracy improved from \textbf{40\% to 95\%} using QLoRA fine-tuning (Phi-2, LoRA rank=16) and few-shot prompting \\[0.4em]
    \textbf{Agent execution (Week 6)} & Gemini 2.5 Flash agent with 9 analytics tools --- \textbf{4/5 scenarios passed (80\%), 6/7 tool selections correct (86\%)} \\[0.4em]
    \textbf{Reproducibility (Week 7)} & ReMindRAG audit: \textbf{17/17 automated tests passed}, 14/14 identified issues resolved \\
    \bottomrule
\end{tabular}
\vspace{0.5em}

\subsection*{What We Learned}

\textbf{The biggest barrier to enterprise trust is not model accuracy --- it is explainability and compliance enforcement.}

Early iterations produced accurate recommendations that managers still could not use, because the system could not cite why a route was chosen or confirm it was physically valid. Once we added the symbolic bridge-constraint veto and the ReMindRAG reasoning trace, confidence in the outputs changed significantly.

The +55 percentage-point accuracy improvement from baseline to adapted model confirmed that domain-specific instruction tuning on logistics jargon and DOT constraint reasoning is not optional --- it is the core product differentiator.

% ─── Section 4 ────────────────────────────────────────────────────────────────
\section{Team \& Commitment}
\textit{Who is on the founding team and why are you working on this venture?}

\vspace{0.3em}
\begin{tabular}{@{}L{3.5cm}L{4.5cm}L{6cm}@{}}
    \toprule
    \textbf{Member} & \textbf{Role} & \textbf{Stack Ownership} \\
    \midrule
    \textbf{Tony Nguyen} & ML / Full-Stack Engineer & Pipeline automation, model adaptation (QLoRA), S3/Snowflake orchestration, reproducibility \\[0.4em]
    \textbf{Daniel Evans} & Data / Backend Engineer & External data engineering (NOAA weather), Streamlit app delivery, evaluation support \\[0.4em]
    \textbf{Joel Vinas} & Data / ML Engineer & Internal data engineering, ingestion scripts, silver-layer preprocessing \\
    \bottomrule
\end{tabular}
\vspace{0.5em}

All three are MS Data Science \& Analytics students with equal equity and full-stack ownership. We are not advisors with slide decks --- we are the engineers who built every component in this system. Responsibilities map directly to code and architecture ownership across the full stack.

\textbf{Why we are working on this:}

We identified the prediction-action gap firsthand while building logistics AI systems for coursework. The constraint-compliance problem is not academic --- it is a live operational risk for every freight company running middle-mile operations. We kept building because the technical path was clear and the operating need was real.

We are committed to this venture because the hardest parts are already built: the data pipeline, the retrieval architecture, the safety veto, the adapted model, and the live application. What remains is go-to-market execution.

% ─── Section 5 ────────────────────────────────────────────────────────────────
\section{RVCC Fit}
\textit{How will RVCC help you move the venture forward?}

\medskip
RVCC fills the precise gaps that technical founders face after building a working product but before achieving first commercial traction.

\subsection*{What We Are Asking For}

\textbf{Mentorship --- go-to-market for logistics operations buyers}\\
We understand the technology. We need guidance on how to sell into operations leadership at carriers and 3PLs: who champions this internally, what procurement looks like at mid-size carriers, and how to structure a paid pilot agreement.

\textbf{Network --- freight operator and ecosystem introductions}\\
A warm introduction to one logistics operator, regional carrier, or Snowflake ecosystem partner would compress our sales cycle significantly. The product is live and demonstrable today.

\textbf{Funding --- runway for hardening and first pilot}\\
Specific near-term needs:
\begin{itemize}[leftmargin=1.5em, itemsep=0.2em]
    \item Complete evaluation hardening and lock benchmark reporting
    \item Scale adaptation dataset beyond 100 examples for broader coverage
    \item Add enterprise access controls (OAuth / RBAC) required for production deployment
    \item Support a first paid pilot engagement
\end{itemize}

\subsection*{Why RVCC Resources Convert Directly to Traction}

\vspace{0.3em}
\begin{tabular}{@{}L{5.5cm}L{8cm}@{}}
    \toprule
    \textbf{Current Asset} & \textbf{RVCC Unlock} \\
    \midrule
    Live deployed application & Pilot-ready with one enterprise introduction \\[0.3em]
    95\% adapted accuracy & Publishable benchmark with mentorship on framing \\[0.3em]
    17/17 reproducibility tests & Engineering credibility for enterprise procurement evaluation \\[0.3em]
    9-tool analytics foundation & Expansion path for enterprise analytics add-on revenue \\
    \bottomrule
\end{tabular}
\vspace{0.5em}

RVCC at this stage would convert proven technical capability into a go-to-market motion --- which is exactly the gap the competition is designed to close.

\vspace{2em}
\begin{center}
    \hrule height 0.8pt
    \vspace{0.5em}
    \textit{EVN HyperLogistics --- RVCC 2026 Submission}
\end{center}

\end{document}
